\subsubsection{CycleGAN Learns to ``Hide'' Information~\citep{chu2017cycleganmastersteganography}}
\label{sec:cyclegan}

Generative Adversarial Networks (GANs; \citealp{goodfellow2014generative}) are generative models that try to learn to produce synthetic data matching the training set's data distribution.
In a GAN, two conjoined models compete in a zero-sum game: a generator seeks to create data matching the training set, while a discriminator attempts to determine if a given data point is from the training-set data distribution or instead was synthetically generated by the generator.
If trained on photorealistic images or 19th-century fiction, the model will learn to replicate the specific styles and structures of those datasets and thus learn to create new instances of photorealistic images or 19th-century fiction.

CycleGAN~\citep{zhu2020unpairedimagetoimagetranslationusing}, a technique used to learn transformations between two distinct image distributions, trains two GANs simultaneously using a ``cycle-consistency'' loss. This loss requires that an image be translatable from its source distribution (e.g., airplane photos) to a target distribution (e.g., Van Gogh paintings) and back again, and afterwards be indistinguishable from its original. However, \citet{chu2017cycleganmastersteganography} observed that CycleGAN often satisfies this constraint by learning to hide detailed information about the source image within the generated target image. % using a nearly imperceptible, high-frequency signal.

For their experiments, \citet{chu2017cycleganmastersteganography} chose aerial photos and maps as their source and target distributions, respectively. CycleGAN was tasked with translating aerial photos into maps and back again, ensuring the final reconstructed images were indistinguishable from the originals.

Aerial photos contain far more complex information than a simplified map; therefore, the model's job is to compress the aerial photo into the map. Because it is physically impossible to cram all the details of a high-resolution photo into the flat colors of a map, the model finds a loophole. Instead of learning the high-level semantic transformation intended---like ``this cluster of pixels represents a building''---the generator learns to hide a full-resolution blueprint of the source image within the generated target image using a nearly imperceptible, high-frequency signal.

To a human observer, the signal is invisible. The generated map looks exactly like a map should, which is why it successfully fools the discriminator, at least initially. For example, a pattern of black dots on a white roof in an original aerial photo was perfectly reconstructed in the final output, even though the corresponding area of the intermediate map appeared as a solid, featureless gray to the naked eye. The model had cached the dot pattern in the pixel noise of the gray patch.

This internal steganography allows the generator to recover the original sample from its transformed counterpart and satisfy the cyclic consistency requirement without the intermediate image generator actually learning the high-level semantic transformation intended.
By viewing this training procedure as generating adversarial examples, \citet{chu2017cycleganmastersteganography} demonstrated that the cyclic consistency loss makes CycleGAN especially vulnerable to adversarial attacks that exploit these side channels.
